\par\Needspace{6\baselineskip}\medskip
\hypertarget{native-115-article086}{}\pdfbookmark[2]{86.}{article-086}
\begin{center}86.\end{center}
Sei $v-\Pi$ die wahre Anomalie, und bezeichnen wir den zugehörigen unbestimmten oder veränderlichen Radiusvektor mit $\rho$. Der Flächeninhalt des Sektors, den der Himmelskörper in der Zeit $t$ überstreicht, ist $=\tfrac12\int\rho\rho\,\mathrm{d}v$, wobei von $v=N$ bis $v=N'$ integriert wird. Nimmt man $k$ in der Bedeutung, die ihm in Art. 6 gegeben wurde, so gilt also $kt\BareRoot{p}=\int\rho\rho\,\mathrm{d}v$. Bezeichnet $\varphi x$ eine beliebige Funktion von $x$, so liefern bekanntlich die von \textsc{Cotes} entwickelten Formeln immer genauere Näherungen für das Integral $\int\varphi x\mathbin{.}\mathrm{d}x$ von $x=u$ bis $x=u+\Delta$. Die Formeln lauten
\[
\begin{gathered}
\begin{aligned}
&\tfrac12\Delta\{\varphi u+\varphi(u+\Delta)\}\\[4pt]
&\tfrac16\Delta\{\varphi u+4\varphi(u+\tfrac12\Delta)+\varphi(u+\Delta)\}\\[4pt]
&\tfrac18\Delta\{\varphi u+3\varphi(u+\tfrac13\Delta)+3\varphi(u+\tfrac23\Delta)+\varphi(u+\Delta)\}
\end{aligned}\\[4pt]
\text{usw.}
\end{gathered}
\]
\noindent Für unseren Zweck genügen die ersten beiden Formeln.

Für unsere Aufgabe liefert die erste Formel also $\int\rho\rho\,\mathrm{d}v=\tfrac12\Delta(rr+r'r')=\dfrac{\Delta rr'}{\cos2\omega}$, wenn man $\dfrac{r'}r=\operatorname{tang}(45^0+\omega)$ setzt. Damit erhält man als ersten Näherungswert von $\BareRoot{p}$ den Ausdruck $=\dfrac{\Delta rr'}{kt\cos2\omega}$, den wir $=3a$ setzen.

Die zweite Formel liefert die genauere Näherung $\int\rho\rho\,\mathrm{d}v=\tfrac16\Delta(rr+r'r'+4RR)$. Dabei bezeichnet $R$ den Radiusvektor, der zur dazwischenliegenden Anomalie $\tfrac12N+\tfrac12N'-\Pi$ gehört. Drückt man nun $p$ nach der in Art. 82 angegebenen Formel durch $r$, $R$, $r'$, $N$, $N+\tfrac12\Delta$, $N+\Delta$ aus, so erhält man
\[
p=\frac{4\sin\tfrac14\Delta^2\sin\tfrac12\Delta}{\left(\dfrac1r+\dfrac1{r'}\right)\sin\tfrac12\Delta-\dfrac1R\sin\Delta}\mathpunct{\raisebox{.45ex}{,}}
\]
\SourceLeaf{116}und daraus
\[
\begin{aligned}
\frac{\cos\tfrac12\Delta}{R}&=\tfrac12\left(\frac1r+\frac1{r'}\right)-\frac{2\sin\tfrac14\Delta^2}{p}\\[4pt]
&=\frac{\cos\omega}{\BareRoot{(rr'\cos2\omega)}}-\frac{2\sin\tfrac14\Delta^2}{p}\mathpunct{\cdot}
\end{aligned}
\]
Setzt man also
\[
\frac{2\sin\tfrac14\Delta^2\BareRoot{(rr'\cos2\omega)}}{\cos\omega}=\delta,
\]
so erhält man
\[
R=\frac{\cos\tfrac12\Delta\mathbin{.}\BareRoot{(rr'\cos2\omega)}}{\cos\omega\mathbin{.}\left(1-\dfrac{\delta}{p}\right)}\mathpunct{\raisebox{.45ex}{,}}
\]
woraus sich der zweite Näherungswert von $\BareRoot{p}$ ergibt
\[
\BareRoot{p}=a+\frac{2a\cos\tfrac12\Delta^2\cos2\omega^2}{\cos\omega^2\left(1-\dfrac{\delta}{p}\right)^2}=a+\frac{\epsilon}{\left(1-\dfrac{\delta}{p}\right)^2}\mathpunct{\raisebox{.45ex}{,}}
\]
wenn man $2a\left(\dfrac{\cos\tfrac12\Delta\cos2\omega}{\cos\omega}\right)^2=\epsilon$ setzt. Schreibt man also $\pi$ für $\BareRoot{p}$, so lässt sich $\pi$ aus der Gleichung $(\pi-a)\left(1-\dfrac{\delta}{\pi\pi}\right)^2=\epsilon$ bestimmen, die nach dem Ausmultiplizieren vom fünften Grad wäre. Setzen wir $\pi=q+\mu$, wobei $q$ ein Näherungswert von $\pi$ und $\mu$ eine sehr kleine Größe ist, deren Quadrat und höhere Potenzen man vernachlässigen kann.

Diese Substitution ergibt
\[
\begin{aligned}
&(q-a)\left(1-\frac{\delta}{qq}\right)^2\\[4pt]
&\quad+\mu\left\{\left(1-\frac{\delta}{qq}\right)^2+\frac{4\delta(q-a)}{q^3}\left(1-\frac{\delta}{qq}\right)\right\}=\epsilon,
\end{aligned}
\]
oder
\[
\mu=\frac{\epsilon q^5-(qq-aq)(qq-\delta)^2}{(qq-\delta)(q^3+3\delta q-4a\delta)}\mathpunct{\raisebox{.45ex}{,}}
\]
und damit
\[
\pi=\frac{\epsilon q^5+(qq-\delta)(aqq+4\delta q-5a\delta)q}{(qq-\delta)(q^3+3\delta q-4a\delta)}\mathpunct{\cdot}
\]
Für unsere Aufgabe liegt bereits ein Näherungswert von $\pi$ vor, nämlich $=3a$. Setzt man diesen Wert in der vorhergehenden Formel für $q$ ein, so erhält man den verbesserten Wert
\[
\pi=\frac{243a^4\epsilon+3a(9aa-\delta)(9aa+7\delta)}{(9aa-\delta)(27aa+5\delta)}\mathpunct{\cdot}
\]
Setzt man also $\dfrac{\delta}{27aa}=\beta$ und $\dfrac{\epsilon}{(1-3\beta)a}=\gamma$, so nimmt die Formel die Gestalt $\pi=\dfrac{a(1+\gamma+21\beta)}{1+5\beta}$ an. Die folgenden fünf Formeln enthalten alle zur Lösung der Aufgabe erforderlichen Rechenschritte.
\[
\begin{aligned}
\mathrm{I.}\qquad&\frac{r'}{r}=\operatorname{tang}(45^0+\omega)\\[6pt]
\mathrm{II.}\qquad&\frac{\Delta rr'}{3kt\cos2\omega}=a
\end{aligned}
\]
\SourceLeaf{117}
\[
\begin{aligned}
\mathrm{III.}\qquad&\frac{2\sin\tfrac14\Delta^2\BareRoot{(rr'\cos2\omega)}}{27aa\cos\omega}=\beta\\[6pt]
\mathrm{IV.}\qquad&\frac{2\cos\tfrac12\Delta^2\cos2\omega^2}{(1-3\beta)\cos\omega^2}=\gamma\\[6pt]
\mathrm{V.}\qquad&\frac{a(1+\gamma+21\beta)}{1+5\beta}=\BareRoot{p}.
\end{aligned}
\]

Begnügt man sich mit etwas geringerer Genauigkeit als der dieser Formeln, so lassen sich noch einfachere Ausdrücke entwickeln. Setzt man $\cos\omega$ und $\cos2\omega=1$ und entwickelt den Wert von $\BareRoot{p}$ in eine Reihe nach Potenzen von $\Delta$, so erhält man nach Vernachlässigung der vierten und höheren Potenzen
\[
\BareRoot{p}=a\left(3-\tfrac12\Delta\Delta+\frac{\Delta\Delta\BareRoot{rr'}}{18aa}\right),
\]
Hier ist $\Delta$ in Teilen des Radius auszudrücken. Setzt man daher $\dfrac{\Delta rr'}{kt}=\BareRoot{p'}$, so ergibt sich
\[
\mathrm{VI.}\qquad p=p'\left(1-\tfrac13\Delta\Delta+\frac{\Delta\Delta\BareRoot{rr'}}{3p'}\right).
\]
Entwickelt man ebenso $\BareRoot{p}$ in eine Reihe nach Potenzen von $\sin\Delta$ und setzt $\dfrac{rr'\sin\Delta}{kt}=\BareRoot{p''}$, so erhält man
\[
\mathrm{VII.}\qquad\BareRoot{p}=\left(1+\frac{\sin\Delta^2\BareRoot{rr'}}{6p''}\right)\BareRoot{p''}
\]
oder
\[
\mathrm{VIII.}\qquad p=p''+\tfrac13\sin\Delta^2\BareRoot{rr'}.
\]
Die Formeln VII und VIII stimmen mit denen überein, die der berühmte \textsc{Euler} in \textit{Theoria motus planetarum et cometarum} [S. 56] angegeben hat. Formel VI stimmt mit derjenigen überein, die in \textit{Recherches et calculs sur la vraie orbite elliptique de la comète de 1769}, S. 80, verwendet wurde.
